% OCEANS 2026 Monterey -- 4-6 pages incl. figures, references excluded.
% Framing: a PIPELINE that turns raw aquarium video into BEHAVIOURAL UNDERSTANDING.
% Long-form source: octopus_behaviour_pipeline_v2.tex (9pp) + PAPER_NOTES.md R1-R35.
% Every number here traces to a committed script; see paper-plan.md for provenance.
\documentclass[conference]{IEEEtran}
\IEEEoverridecommandlockouts
\usepackage{cite}
\usepackage{amsmath,amssymb,amsfonts}
\usepackage{graphicx}
\usepackage{textcomp}
\usepackage{xcolor}
\usepackage{booktabs}
\usepackage{multirow}
\renewcommand{\topfraction}{0.92}
\renewcommand{\bottomfraction}{0.7}
\renewcommand{\textfraction}{0.06}
\renewcommand{\floatpagefraction}{0.85}
\begin{document}

\title{From Footage to Ethogram: A Deployable Pipeline for Continuous
Behavioural Monitoring of a Captive Octopus}

\author{
\IEEEauthorblockN{
Siddharth Raj\IEEEauthorrefmark{1},
Harald Burgsteiner\IEEEauthorrefmark{2},
Wolfgang Slany\IEEEauthorrefmark{3}
}
\IEEEauthorblockA{\IEEEauthorrefmark{1}
Research and Development, Arishna IoT Solutions, Bangalore, India}
\IEEEauthorblockA{\IEEEauthorrefmark{2}
Institute for Digital Media Education, University of Teacher Education Styria, Graz, Austria}
\IEEEauthorblockA{\IEEEauthorrefmark{3}
Institute of Software Engineering and Artificial Intelligence, Graz University of Technology, Graz, Austria}
}
\maketitle

\begin{abstract}
Continuous behavioural monitoring is what an aquarium actually wants from its
cameras, and it is out of reach: manual ethogram scoring does not scale, and
prompting a large vision--language model (VLM) on every clip of a single
animal's footage costs hundreds of hours of API time and cannot run on site. We
present a pipeline that turns raw aquarium video into a behavioural time series
using a cascade of small local models, with a $235$B-parameter VLM employed only
as an offline teacher. On $892$\,h of footage from one adult
\emph{Octopus vulgaris}, cheap gates discard $60\%$ of videos without a full
decode and only ${\approx}60$ of the $892$\,h is ever decoded; of the clips that
survive, $42.7\%$ still contain no animal, which is why the expensive model is
spent last rather than first. Each gate is placed by a measurement on
human-labelled data rather than by intuition, and in each case the measurement
changed the design.
The distilled students reach $75.4\%$ accuracy and $0.665$ macro-F1 on a
six-class ethogram, $0.642$ mask IoU with no prompt at inference from
$3.2$\,M parameters, and captioning that improves embedding similarity from
$0.702$ to $0.834$ over its own base model. Applied to the corpus, the pipeline
recovers a circadian activity profile peaking at $45\%$ visible activity at
$17{:}00$ against $1$--$5\%$ overnight, and shows that human presence nearly
doubles movement. We release $5{,}222$ VLM-labelled clips with their full vote
distributions, ${\approx}25{,}000$ captions, $969$ human labels and masks, and
three frozen video-level benchmarks.
\end{abstract}

\begin{IEEEkeywords}
animal behaviour, cephalopods, ethogram, knowledge distillation, video
segmentation, vision--language models
\end{IEEEkeywords}

\section{Introduction}
An aquarium camera produces footage continuously; turning it into behavioural
knowledge does not scale the same way. Manual ethogram scoring is the accepted
standard and is bounded by expert hours. The obvious modern substitute---prompt
a large VLM on every clip---fails on arithmetic rather than on accuracy. The
$892$\,h of colour-camera footage we hold for a single animal is
${\approx}160{,}000$ twenty-second clips; at the throughput we measured against
a hosted $235$B model ($13$--$36$ clips/min sustained), a five-pass labelling
run is $372$--$1{,}030$\,h of wall-clock, and it requires an API the aquarium
does not have. The server holds roughly three times more footage again, for four
further animals.

The waste is also structural, not merely large. Cephalopods spend most of their
time hidden or motionless: after the extraction gates described below,
$42.7\%$ of surviving clips still contain no visible animal and a further
$16.7\%$ show it stationary. Prompting a frontier model uniformly spends most of
its budget confirming that nothing is happening.

We therefore treat the VLM as an \emph{offline teacher} and build a cascade of
small models that runs locally. The contribution is not the cascade as an
architecture---teacher--student distillation is standard---but that every gate in
it is placed by a measurement on human-labelled data, and that several of those
measurements contradicted the intuitive choice. We then show the pipeline
produces behavioural findings, which is the point of building it, and release the
labels and benchmarks so the measurements can be checked.

\section{Related Work}
Our own prior work established the recording setup and the framing of
inter-species communication with octopodes~\cite{disc_o}. For perception we use
CLIP~\cite{clip} features, Grounding DINO~\cite{groundingdino} and
SAM2~\cite{sam2} as annotators, LoRA~\cite{lora} for adaptation,
MobileNetV3~\cite{mobilenetv3} as a deployable backbone, and
Qwen3-VL~\cite{qwenvl} as the teacher.

\textbf{Foundation models as teachers, not detectors.} Open-vocabulary detection
is attractive here because no octopus detector exists off the shelf. In practice
OWLv2~\cite{owlv2} reported an octopus in $231$ of $232$ frames at its default
threshold. Its scores are not uninformative---they rank the two classes at AUC
$0.759$ against human adjudication---but the distributions overlap so heavily
that no threshold is usable: at its own optimum it still passes $30\%$ of the
frames we wanted removed while discarding $32\%$ of the real animals. A weak
ranking signal with no viable operating point is not a detector, which is the
case for distillation rather than prompting.

\textbf{Octopus segmentation.} The closest work is
HideAndSeg~\cite{hideandseg}, a concurrent annotation-light pipeline: a human
annotates initial frames to prompt SAM2, the resulting masks train a YOLOv11
detector, and that detector auto-prompts SAM2 over the remaining footage. Its
central insight---re-detecting per frame so a contaminated propagation track can
recover---is correct, and our seed-confidence gate serves the same purpose. Its
field coverage substantially exceeds ours: $148$ videos and $366{,}514$ frames of
wild juvenile \emph{O.\ insularis} across three Brazilian sites, against our
single individual in one tank. We compare against it in
Table~\ref{tab:sidebyside}, where the condition rows matter more than the metric
row. Notably, its stated motivation---the absence of publicly available
annotated octopus video---holds after publication, since it releases neither data
nor code; our release is a partial answer to it.

\section{Pipeline and the Measurements That Set It}
\label{sec:pipeline}
The cascade runs cheapest-first: probe, motion, a trained presence probe, then
the VLM (Fig.~\ref{fig:pipeline}). Table~\ref{tab:gates} reports how much
footage survives each stage and, for every gate, the measurement that placed it
alongside the decision that measurement reversed. In each case the measurement
changed the design we would have chosen without it.

\begin{figure}[t]\centering
  \includegraphics[width=0.47\textwidth]{assets/pipeline_behaviour.pdf}
  \caption{The pipeline. Cheap local gates run first; the teacher VLM is reached
  only by clips that survive them, and is used offline to train the students that
  are actually deployed.}
  \label{fig:pipeline}
\end{figure}

\begin{table}[t]\centering
\caption{The Cascade: Survival Rates and Gate Placement}
\label{tab:gates}
\small
\begin{tabular}{@{}p{0.26\columnwidth}p{0.66\columnwidth}@{}}
\toprule
\multicolumn{2}{@{}l}{\textbf{Funnel} (1{,}769 videos, $892.5$\,h)} \\
\midrule
probe & $1{,}054$ videos ($60\%$) discarded on ${\approx}10$ seek frames, with no full decode \\
decoded at all & $\mathbf{{\approx}60}$\,h \\
clips to VLM & $5{,}222$ \\
\quad no animal & $42.7\%$ \\
\quad stationary & $16.7\%$ \\
\quad active & $\mathbf{40.6\%}$ \\
\midrule
\multicolumn{2}{@{}l}{\textbf{Gate placement}} \\
\midrule
presence &
the VLM returns \emph{not present} on $63\%$ of clips the trained detector had
marked verified. Prevented trusting the detector. \\
label confidence &
human agreement tracks the five-vote margin: $0.73$ unanimous, $0.86$ at
four-of-five, $0.43$ at $\leq$three-of-five. Prevented arg-max labels teaching
false confidence; targets are the vote distribution. \\
motion &
per-video \emph{normalised} motion passes static scenes---a flickering lamp
normalises up to ``motion''. Absolute changed-pixel fraction used instead. \\
\bottomrule
\end{tabular}
\end{table}

\textbf{Labelling with five disjoint views.} Each surviving clip is sent to the
teacher five times. Pass $p$ offsets its uniform frame grid by
$(p-1)\,\mathrm{step}/5$, so the five samplings are disjoint and their union
tiles every dense frame; a clip therefore receives five independent presence
votes, ethogram votes and captions. This is not redundancy for its own sake: the
majority vote differs from a single pass on $705$ clips ($13.5\%$), $36.2\%$ of
clips are split rather than unanimous, and---per Table~\ref{tab:gates}---the
margin predicts human agreement, which is what licenses training on soft targets.

\section{Local Models}
\label{sec:models}
Three students consume the teacher's output. All are trained with splits
\emph{by source video}; a held-out clip from a training video is not held out in
any useful sense.

\textbf{Ethogram classifier.} Six classes in one head---\emph{no octopus},
resting, exploration/manipulation, locomotion, reaching out of water,
human/enrichment interaction---because that is how it is deployed: a clip
arrives and the question is what is happening, with ``no animal'' a legitimate
answer. Inputs are frozen per-frame features plus two changed-pixel motion
channels; the head is a $0.6$\,M-parameter MLP over pooled statistics. A
previous attempt at this task collapsed onto the majority classes ($45\%$
accuracy, per-class F1 ${\approx}0$). Table~\ref{tab:etho} shows the current
model and the one change that mattered.

\begin{table}[t]\centering
\caption{Ethogram Classification Results}
\label{tab:etho}
\small
\begin{tabular}{@{}lcc@{}}
\toprule
model & macro-F1 & accuracy \\
\midrule
majority class                              & $0.100$ & $43.1\%$ \\
pooled CLIP $\rightarrow$ linear            & $0.400$ & -- \\
CLIP $\rightarrow$ BiGRU                    & $0.530$ & -- \\
three self-supervised backbones, ensembled  & $0.617$ & $71.4\%$ \\
\textbf{$+$ mask-guided crop (five members)}& $\mathbf{0.665}$ & $\mathbf{75.4\%}$ \\
\midrule
\multicolumn{3}{@{}l}{\emph{per class} (precision / recall)} \\
no octopus                    & \multicolumn{2}{l}{$0.93$ / $0.90$} \\
resting                       & \multicolumn{2}{l}{$0.71$ / $0.60$} \\
exploration/manipulation      & \multicolumn{2}{l}{$0.66$ / $0.67$} \\
locomotion                    & \multicolumn{2}{l}{$0.64$ / $0.73$} \\
reaching out of water         & \multicolumn{2}{l}{$0.57$ / $0.75$} \\
human/enrichment interaction  & \multicolumn{2}{l}{$0.44$ / $0.43$} \\
\bottomrule
\end{tabular}
\end{table}

The gain came from the \emph{representation}, not the head. Replacing CLIP with a
self-supervised image backbone or a video backbone moved test macro-F1 by
$+0.047$ and $+0.059$; the worst configuration of either beat the best CLIP
configuration, with no overlap in range. Head capacity, by contrast, did nothing:
sweeping hidden width $256$--$1024$ moved validation by $+0.0017$ against a seed
standard deviation of $0.0040$, and the widest head was the worst on both
backbones.

A second representation effect isolates \emph{animal pixel resolution} as the
operative variable. The median animal occupies a bounding box $0.216$ of the
frame across, so a backbone at $224^2$ sees it at ${\approx}48\times48$ pixels.
Three preprocessing variants separate the causes: the default centre crop retains
$43\%$ of the frame but preserves animal scale ($0.578$ validation);
letterboxing retains all of it but shrinks the animal ($0.574$); cropping to the
segmentation mask magnifies the animal ${\approx}3\times$ and reaches $0.618$.
Field of view and animal scale trade off almost exactly---only magnification
helps.

\textbf{Segmentation.} A $3.2$\,M-parameter LR-ASPP student distilled from
Grounding DINO$+$SAM2 reaches IoU $0.642$ mean / $0.719$ median against $122$
human-drawn masks on five videos held out of all training, with mean area error
${\approx}1\%$. It does not meet the $0.85$ IoU we set out to reach; we report it
because the shortfall is informative. Everything downstream---presence, body
area as a posture proxy, motion inside the mask---is area-based, and area is
accurate to $1\%$ while boundary quality is what IoU penalises. Two levers
produced the gain, both mundane and both measured leak-free: raising input
resolution from $256^2$ to $512^2$ (at $256^2$ the median animal is
${\approx}40\times40$ pixels and an arm is $1$--$2$ pixels wide, i.e.\
unrepresentable) and replacing symmetric Dice$+$BCE with Focal--Tversky, which
penalises the under-segmentation the symmetric loss ignores.

\begin{figure}[t]\centering
  \includegraphics[width=0.42\textwidth]{assets/segmentation_demo.pdf}
  \caption{Mask student output. Per-frame prompting alone bleeds into bright
  metal tools under infrared and into background clutter on colour cameras;
  confidence-seeded temporal propagation with a largest-blob constraint removes
  both, and correctly rejects the reflection camera.}
  \label{fig:seg}
\end{figure}

\textbf{Captioning.} A $2$B VLM adapted with LoRA on teacher captions improves
held-out embedding similarity from $0.702$ to $0.834$ and ROUGE-L from $0.269$ to
$0.455$ over its own base model. Quantised to four bits it is $1.7$\,GB and
captions a clip in ${\approx}3$\,s on a $16$\,GB laptop with no GPU.

\textbf{What did not work.} Recorded because each is a natural next thing to try.
Feature-space augmentation (mixup, Gaussian noise) is \emph{negative} and
monotonically worse with strength; the model is not overfitting-limited, so a
regulariser only corrupts frozen features. Class rebalancing hurt the class it
was meant to protect---removing it improved the weakest class from F1 $0.37$ to
$0.44$. And our own hypothesis that teacher-label quality was the segmentation
ceiling was wrong: the evidence for it turned out to be train leakage, and on a
clean holdout, mask IoU is flat across teacher quality ($0.494$, $0.508$,
$0.506$).

\begin{table}[t]\centering
\caption{Distilled Students vs.\ Zero-Shot Teachers}
\label{tab:zeroshot}
\small
\begin{tabular}{@{}llc@{}}
\toprule
task & system & metric \\
\midrule
\multirow{2}{*}{mask} & GD$+$SAM2, zero-shot, per frame & IoU $0.374$ \\
 & \textbf{ours}, $3.2$\,M, no prompt & IoU $\mathbf{0.642}$ \\
\multicolumn{3}{@{}l}{\quad paired $\Delta=-0.268$, CI$_{95}$ $[-0.313,-0.136]$} \\
\midrule
\multirow{2}{*}{presence} & zero-shot CLIP, $5{+}5$ prompts & AUC $0.450$ \\
 & \textbf{ours}, trained probe & AUC $\mathbf{0.745}$ \\
\multicolumn{3}{@{}l}{\quad paired $\Delta=+0.295$, CI$_{95}$ $[+0.069,+0.528]$} \\
\midrule
\multirow{2}{*}{caption} & base $2$B VLM & emb-sim $0.702$ \\
 & \textbf{ours}, LoRA & emb-sim $\mathbf{0.834}$ \\
\bottomrule
\end{tabular}
\end{table}

Table~\ref{tab:zeroshot} sets the students against the zero-shot models they were
distilled from, on identical held-out human labels; paired differences are
bootstrapped clustered by source video and both intervals exclude zero. The two
presence arms share the frozen backbone \emph{and} its preprocessing, so the gap
isolates the trained probe rather than the features. One caveat travels with the mask row: on the
$21$ of $122$ frames where Grounding DINO clears its own confidence gate the
teacher \emph{wins}, $0.726$ against $0.657$. The teacher is high-precision and
low-recall---it finds nothing at all on a quarter of frames---and distillation
converted a sparse, high-quality signal into dense coverage. That is a better
argument for distilling than a flat win would have been.

\begin{table}[t]\centering
\caption{Comparison with HideAndSeg~\cite{hideandseg}}
\label{tab:sidebyside}
\small
\begin{tabular}{@{}lccc@{}}
\toprule
 & \cite{hideandseg} & our teacher & \textbf{our student} \\
\midrule
mask IoU              & $0.938$             & $0.726^{a}$ & $\mathbf{0.642}$ \\
reference             & init.\ frame        & human masks & human masks \\
prompt at inference   & auto               & gated       & \textbf{none} \\
presence              & given               & inferred    & \textbf{inferred} \\
params                & ${\approx}249$\,M   & ${\approx}230$\,M & $\mathbf{3.2}$\,M \\
data/code released    & no                  & --          & \textbf{yes} \\
setting               & wild, 3 sites       & 1 tank      & 1 tank \\
scale                 & $148$ vid.          & --          & $4{,}665$ clips \\
\bottomrule
\end{tabular}
\end{table}

Table~\ref{tab:sidebyside} places us beside that work. Reported mask quality
tracks how much prompting is assumed at inference, and every other row differs, so
the IoU row is not a ranking. We claim the unprompted, presence-inferring and
${\approx}78\times$-smaller rows, and the release; we do not claim the IoU row,
since neither method has been run on the other's data and the final two rows
favour them.

\section{What the Pipeline Reveals}
\label{sec:findings}
Running the pipeline over the corpus yields a behavioural record, which is the
reason to build it. Across $3{,}083$ clips with a visible animal, the activity
budget is $41\%$ exploration/manipulation, $33\%$ resting, $14\%$
human/enrichment interaction, $9\%$ reaching out of water and $3\%$ locomotion.

\begin{figure}[t]\centering
  \includegraphics[width=0.48\textwidth]{assets/behaviour_findings.pdf}
  \caption{Behaviour recovered from the corpus: activity budget, and
  exposure-normalised visible activity by hour of day. Normalising by the number
  of windows examined per hour is essential---without it the profile tracks
  recording coverage rather than the animal.}
  \label{fig:findings}
\end{figure}

\textbf{Circadian structure.} Normalising visible activity by the number of
windows actually examined in each hour---without which the profile simply tracks
recording coverage---gives a strongly diurnal animal: $1$--$5\%$ visible activity
overnight rising to a $45\%$ peak at $17{:}00$, with a plateau from $13{:}00$ to
$19{:}00$ and a smaller rise at dawn.

\textbf{Response to human presence.} Clips the extractor labels as containing a
human show nearly double the absolute motion of baseline clips ($0.095$ against
$0.045$) and a higher score on a transparent arousal rubric ($0.68$ against
$0.46$). On the colour cameras, the most common body colour at baseline
(dark red-brown, $16\%$ of clips) is markedly less frequent during human
interaction ($6\%$).

Two caveats bound these. The tank contains permanent enrichment objects, so a
context label of ``enrichment object present'' fires on ${\approx}66\%$ of clips
and means only that---not active enrichment; the clean contrast is human present
against absent. And the upstream presence gate is imperfect, so absolute levels
will shift with a better detector while the \emph{rate contrasts} above, which
are ratios within the same gate, are robust to it.

\section{Released Data and Benchmarks}
\label{sec:data}
No public annotated octopus behaviour corpus exists. We release:
$5{,}222$ clips with five-pass teacher labels \emph{including the full vote
distribution} rather than only the majority; ${\approx}25{,}000$ captions
($5{,}160$ clips $\times$ five disjoint frame samplings), which additionally
permits measuring how a VLM's description of one behaviour varies with which
frames it happens to see; a frozen six-class ethogram dataset of $4{,}665$ clips
over $204$ videos with video-level splits and per-clip reliability weights;
$456$ human behaviour labels; $513$ human masks; and $3{,}205$ structured
behavioural records of nine fields each. Three benchmarks are frozen: $122$ human
masks on five held-out videos plus $19$ empty-tank negatives; $50$ human
arm-keypoint frames over $20$ videos; and a reflection-rejection set.

Four protocol rules are part of the release, because the pipeline broke each of
them at least once during development: splits are by source video, never by clip;
frozen sets are never regenerated to suit a result; negatives of different kinds
are never pooled; and holdout videos are excluded from \emph{every} training
source, not merely the final stage.

\section{Toward Limb-Level Ethograms}
The mask student also supports a finer representation than a class label. We
convert a silhouette into an anatomical graph---mantle, head and eight
arms---by thinning and shortest-path tracing, tracking it across frames with an
optical-flow prior and marking every node \emph{detected}, \emph{fitted} or
\emph{occluded} so that kinematics use only evidence-backed samples. On $50$
human-annotated frames the arm-tip F1 is $0.539$ (precision $0.72$) at $3.68$
arms recovered per frame; seeding propagation from the best-resolved frame rather
than the first was worth $+3.67$ arms. Speeds are in pixels per second in crop
space, uncalibrated, and arm counts are silhouette-limited.

\begin{figure}[t]\centering
  \includegraphics[width=0.44\textwidth]{assets/skeleton_example.pdf}
  \caption{Silhouette to anatomical graph: mantle, head and eight arms, with each
  node marked detected, fitted or occluded so that kinematics are computed only
  from evidence-backed samples.}
  \label{fig:skel}
\end{figure}

This is released as a capability rather than a result. Its purpose is the obvious
next question: a six-class sheet cannot distinguish exploration with one arm from
whole-body exploration, and per-limb usage statistics would make
\emph{limb-aware} ethograms possible. That is the direction we intend to take.

\section{Limitations}
\label{sec:limits}
\textbf{One animal, one tank.} Every measurement is from a single adult
\emph{O.\ vulgaris}. Footage for four further animals, including $155$ days of a
second individual, is available and unprocessed; cross-animal generalisation is
untested, and this is the paper's principal limitation.

\textbf{Agreement, not accuracy.} All $456$ human behaviour labels were collected
with the model's suggestion visible, so they measure agreement with the model, not
independent accuracy. On the subset where the two disagree, the human sides with
the teacher over the student roughly $3{:}2$. A blind round on the reserved test
videos is specified but not yet run, and until it is, the ethogram figures should
be read as teacher reproduction.

\textbf{Selection on a small validation set.} Roughly six design choices were
made on $35$ validation videos, and validation and test disagreed on three of
them. Reported margins are softer than their seed variance suggests; we report
the validation-selected configuration throughout rather than the best test score.

\textbf{Infrared.} The colour-trained mask student is unusable on the infrared
camera, which is $35\%$ of the corpus, so mask-derived quantities are available
for colour cameras only.

\section{Conclusion}
Continuous behavioural monitoring of a captive cephalopod is achievable if the
expensive model is used to teach rather than to run. Placing each gate in the
cascade by measurement rather than intuition mattered more than any single model
choice: every gate ended up somewhere other than the obvious default, and the
largest single improvement to behaviour classification came from changing the
frozen representation and how much of the animal it sees, not from any change to
the classifier. The resulting local models recover circadian structure and a
clear response to human presence from footage that no one scored by hand. Data,
benchmarks and code are released; the immediate next steps are a blind human
round on the reserved videos and the second animal.

\section*{Ethical Statement and Data Availability}
All footage is observational, from cameras already installed for husbandry, with
no intervention or manipulation of the animal. Datasets, frozen benchmarks,
weights and an extended technical report are available at the project
repository.

\begin{thebibliography}{00}
\bibitem{disc_o}
K.\ Patel, A.\ Sharma, H.\ Kumar, P.\ Kurundodi, H.\ Burgsteiner, and W.\ Slany,
``A Novel Framework for Decoding Inter-Species Communication with Octopodes
(DISC-O),'' in \textit{OCEANS 2025}, pp.\ 1--9, Sep.\ 2025,
doi: 10.23919/OCEANS59106.2025.11244956.
\bibitem{clip}
A.\ Radford \textit{et al.}, ``Learning Transferable Visual Models from Natural
Language Supervision,'' in \textit{ICML}, 2021.
\bibitem{groundingdino}
S.\ Liu \textit{et al.}, ``Grounding DINO: Marrying DINO with Grounded Pre-Training
for Open-Set Object Detection,'' \textit{arXiv:2303.05499}, 2023.
\bibitem{sam2}
N.\ Ravi \textit{et al.}, ``SAM 2: Segment Anything in Images and Videos,''
\textit{arXiv:2408.00714}, 2024.
\bibitem{lora}
E.\ J.\ Hu \textit{et al.}, ``LoRA: Low-Rank Adaptation of Large Language Models,''
in \textit{ICLR}, 2022.
\bibitem{mobilenetv3}
A.\ Howard \textit{et al.}, ``Searching for MobileNetV3,'' in \textit{ICCV}, 2019.
\bibitem{qwenvl}
Qwen Team, ``Qwen3-VL,'' 2025.
\bibitem{owlv2}
M.\ Minderer, A.\ Gritsenko, and N.\ Houlsby, ``Scaling Open-Vocabulary Object
Detection,'' in \textit{NeurIPS}, 2023.
\bibitem{hideandseg}
A.\ de Aguiar, M.\ P.\ Andrade, C.\ M.\ D.\ Santos, and J.\ P.\ Gois,
``HideAndSeg: an AI-based tool with automated prompting for octopus segmentation
in natural habitats,'' \textit{arXiv:2511.04426}, Nov.\ 2025.
\end{thebibliography}

\end{document}
